% id: PHY-0038
% title: การกระจัดมากที่สุดและกราฟความเร่งจากกราฟความเร็ว
% camp: อื่น ๆ
% topic: การเคลื่อนที่แนวตรง
% year: 2552
% tags: IJSO, รอบที่ 2, การกระจัด, ความเร่ง, กราฟความเร็ว–เวลา
% origin: IJSO รอบที่ 2 2552 · ข้อ 38 ในภาพที่ผู้ใช้ส่ง
% review: ถอดจากภาพแล้ว · รอตรวจทาน
% diagram: PHY-0038.svg
% scan: sources/PHY-0038-original.png
\documentclass[12pt]{article}
\usepackage{amsmath,amssymb,fontspec,polyglossia,tikz}
\setdefaultlanguage{thai}
\setmainfont{Noto Serif Thai}
\usetikzlibrary{arrows.meta}
\begin{document}
\textbf{ข้อ 38.} (IJSO รอบที่ 2 2552)
ข้างล่างเป็นกราฟของความเร็วเทียบกับเวลาของวัตถุหนึ่งซึ่งกำลังเคลื่อนที่ในแนวเส้นตรง

ก. การกระจัดของวัตถุจากตำแหน่งตั้งต้นมีขนาดมากที่สุดเท่าใด

ข. จงเขียนกราฟของความเร่งของวัตถุเทียบเท่ากับเวลา

\begin{tikzpicture}[x=1.75399cm,y=0.51111cm,>=Stealth]
\draw[->,line width=.7pt] (-0.18,0) -- (3.55,0);
\draw[->,line width=.7pt] (0,-2.8) -- (0,5);
\draw[line width=.7pt] (1,-0.1304) -- (1,0.1304);
\node[anchor=center] at (1,-0.5217) {1};
\draw[line width=.7pt] (2,-0.1304) -- (2,0.1304);
\node[anchor=center] at (2,-0.5217) {2};
\draw[line width=.7pt] (3,-0.1304) -- (3,0.1304);
\node[anchor=center] at (3,-0.5217) {3};
\draw[line width=.7pt] (-0.038,-2) -- (0,-2);
\node[anchor=east] at (-0.0912,-2) {$-2$};
\draw[line width=.7pt] (-0.038,0) -- (0,0);
\node[anchor=east] at (-0.0912,0) {0};
\draw[line width=.7pt] (-0.038,2) -- (0,2);
\node[anchor=east] at (-0.0912,2) {2};
\draw[line width=.7pt] (-0.038,4) -- (0,4);
\node[anchor=east] at (-0.0912,4) {4};
\draw[line width=.7pt] (0.5,-0.1304) -- (0.5,0);
\draw[line width=.7pt] (1.5,-0.1304) -- (1.5,0);
\draw[line width=.7pt] (2.5,-0.1304) -- (2.5,0);
\draw[line width=.7pt] (-0.038,-1) -- (0,-1);
\draw[line width=.7pt] (-0.038,1) -- (0,1);
\draw[line width=.7pt] (-0.038,3) -- (0,3);
\node[anchor=west] at (3.6108,-0.0522) {เวลา (s)};
\node[anchor=west] at (-0.0608,5.4174) {ความเร็ว (m/s)};
\draw[line width=.7pt] (0,2) -- (1,4) -- (1.5,4) -- (3,-2);
\end{tikzpicture}
\end{document}
